\chapter*{Abstract}
\addcontentsline{toc}{chapter}{Abstract}

Travel planning is a multi-stage information task in which users must identify
destinations, compare alternatives, satisfy personal constraints, and combine
information from heterogeneous sources. General-purpose Large Language Models
(LLMs) can produce fluent recommendations, but their answers may be outdated,
unsupported, or insufficiently adapted to the user's preferences. This thesis
develops a Personalized Travel Guide Assistant for Vietnam using
Retrieval-Augmented Generation (RAG), persistent user memory, temporary
conversation state, and requirement-aware evidence checking.

The proposed system constructs a tourism knowledge base from webpages and
documents using content cleaning, structure-aware chunking, LLM-assisted
metadata extraction, sentence embeddings, and PostgreSQL with vector-search
support. During interaction, the system rewrites conversational follow-ups,
extracts intent, operation, constraints, and semantic retrieval facets, and
creates bounded retrieval tasks for complex requests. Dense semantic retrieval
and BM25 lexical retrieval are combined using Reciprocal Rank Fusion, followed
by CrossEncoder reranking. An evidence checker measures whether the selected
chunks cover each factual requirement. Missing evidence triggers bounded
recovery retrieval, while an answer-readiness gate determines whether the
system should produce a complete answer, a qualified partial answer, or a
knowledge-gap response.

The evaluation framework separates query understanding, retrieval quality, and
final-answer quality. It uses a controlled multi-turn dataset consisting of five
synthetic users, two conversations per user, and ten queries per conversation,
for 100 query turns. The principal metrics are intent accuracy, operation
accuracy, query-constraint F1, retrieval-facet F1, nDCG@5, Precision@5,
correctness, faithfulness, personalization adherence, and completeness. An
independent LLM provides initial judgments, followed by human validation of
reference labels and selected outputs. The experiment additionally records
latency, token usage, estimated cost, model versions, prompt hashes, and a
corpus fingerprint.

The work contributes an inspectable architecture for personalized,
evidence-aware travel assistance and a reproducible evaluation design that
connects final-answer behavior to query parsing, memory, retrieval, coverage,
and corpus limitations.

\textbf{Keywords:} retrieval-augmented generation, personalized travel
assistant, conversational recommendation, hybrid retrieval, user memory,
evidence coverage, large language model evaluation.

